\section{Grundlagen, Modell und Daten}

\subsection{Mechanistische Interpretierbarkeit}

Die meiste Forschung zu Sprachmodellen beschränkt sich auf die Beziehung zwischen Eingabe und Ausgabe. Man beobachtet, was ein Modell tut, ohne zu wissen, wie es dazu kommt. Die \gls{mechanisticinterpretability} geht einen Schritt weiter und öffnet die sprichwörtliche Black Box. Ein Bild, das den Unterschied gut greifbar macht, ist der zwischen dem Fahren eines Autos und dem Verstehen seines Motors. Erst das Verständnis der inneren Abläufe erlaubt es, Fehler vorherzusagen und gezielt einzugreifen \cite{nanda2023comprehensive}.

\subsection{Residual Stream und fünf Werkzeuge}

Moderne \glspl{transformer} bestehen aus einem Stapel von Schichten, die alle auf einen gemeinsamen Informationskanal zugreifen, den \emph{Residual Stream}. Jede Schicht liest daraus, berechnet einen Beitrag und addiert ihn wieder hinzu, ohne den bisherigen Inhalt zu löschen \cite{elhage2021mathematical}. Man kann sich das wie ein Notizblatt vorstellen, auf das jede Schicht etwas ergänzt, statt es zu überschreiben; am Ende wird aus diesem Notizblatt die Vorhersage abgelesen.

Die Arbeit nutzt darauf aufbauend fünf Werkzeuge, die von beschreibenden hin zu kausalen Aussagen führen. Die \textbf{Residual-Stream-Analyse} und das \textbf{lineare Probing} prüfen, ob Sentiment in den \glspl{activation} vorhanden und linear dekodierbar ist, also ob schon eine einfache Trennlinie reicht, um positiv von negativ zu unterscheiden. Der \textbf{Logit Lens} projiziert Zwischenrepräsentationen einzelner Schichten in den Vokabularraum \cite{nostalgebraist2020logit} und zeigt so, ab wann das Modell eine Information nutzt. Die \textbf{Attention-Analyse} zeigt, welche Köpfe auf welche Token zugreifen, also wohin Information \emph{fließt}. Nur das \textbf{Activation Patching} weist nach, dass eine Komponente kausal \emph{benötigt} wird, indem einzelne \glspl{activation} gezielt ersetzt und die Wirkung auf die Vorhersage gemessen wird \cite{meng2022locating, goldowskydill2023localizing}. Die ersten vier Werkzeuge liefern Korrelationen, nur das letzte liefert einen echten kausalen Nachweis.

Für die Sentiment-Bewertung wird durchgehend das Opinion-Lexikon von Hu und Liu verwendet, das positive und negative Wörter wie „good“/„bad“ oder „excellent“/„terrible“ enthält \cite{hu2004mining}.

\subsection{Das Modell Pythia-410M}

\begin{table}[H]
  \centering
  \caption{Zentrale Architekturwerte von Pythia-410M.}
  \label{tab:arch}
  \begin{tabular}{ll}
    \toprule
    Eigenschaft & Wert \\
    \midrule
    Architekturtyp & Decoder-Only (GPT-NeoX) \\
    Vokabulargröße & 50.304 Token \\
    Hidden Size & 1.024 \\
    Transformer-Schichten & 24 \\
    \glspl{attentionhead} je Schicht & 16 \\
    Positionskodierung & Rotary Positional Embeddings (RoPE) \\
    \bottomrule
  \end{tabular}
\end{table}

Pythia-410M ist ein offenes, vollständig dokumentiertes Decoder-Only-\gls{transformer}-Modell (GPT-NeoX-Architektur) von EleutherAI mit rund 405 Millionen Parametern \cite{biderman2023pythia}. Diese verteilen sich auf vier Hauptkomponenten. Embedding- und Unembedding-Matrix tragen je rund 12,7~Prozent, die \gls{attention}-Blöcke über alle 24 Schichten rund 24,9~Prozent, und die Feed-Forward-Netzwerke vereinen mit rund 49,7~Prozent fast die Hälfte der Parameter auf sich. Der größte Teil der Modellkapazität liegt also nicht in der Attention selbst, sondern in den Feed-Forward-Blöcken. Tabelle~\ref{tab:arch} fasst die zentralen Architekturwerte zusammen.


\subsection{Datensätze}

Für die verschiedenen Analysen dieser Arbeit werden drei sich ergänzende Datensätze verwendet. \par

Das \textbf{Hu-\&-Liu-Opinion-Lexikon} liefert 2.254 Ein-Token-Wörter (865 positiv, 1.389 negativ), die vom Tokenizer von Pythia-410M jeweils als genau ein Token dargestellt werden. Das ursprüngliche Lexikon umfasst rund 6.800 sentimenttragende Wörter. Die Filterung erfolgt, indem jedes Wort mit führendem Leerzeichen tokenisiert und nur dann übernommen wird, wenn genau eine Token-ID erzeugt wird. Dadurch kann jedem Wort eindeutig ein Embedding-Vektor sowie ein einzelner Logit-Wert zugeordnet werden. Das Lexikon dient in dieser Arbeit zur Identifikation sentimentbezogener Tokens sowie zur Berechnung von Sentiment-Scores und Logit-Differenzen über die Transformer-Schichten hinweg \cite{hu2004mining}. \par

Das \textbf{Sentiment Challenge Dataset} enthält 836 binär gelabelte Sätze, die gezielt sprachliche Herausforderungen wie Negation, Ironie, Idiome, Kompositionalität oder gemischtes Sentiment abdecken. Der Datensatz wird verwendet, um das Verhalten des Modells unter kontrollierten sentimentbezogenen Bedingungen zu untersuchen und dient als Grundlage für die qualitativen Analysen in Kapitel 3 \cite{barnes2019sentiment}. \par

Die \textbf{Counterfactually Augmented Data (CAD)} bestehen aus 1.707 Paaren positiver und negativer Texte, die sich nur in wenigen sentimenttragenden Wörtern unterscheiden. Da beide Varianten nahezu identischen Kontext besitzen, eignen sie sich besonders gut zur isolierten Untersuchung von Sentiment-Repräsentationen. In dieser Arbeit werden die CAD-Paare für die aggregierten Logit-Lens-, Attention-Head-, Linear-Probing- und Activation-Patching-Analysen verwendet \cite{kaushik2020learning}.
